\documentclass[12pt]{article}

\usepackage[utf8]{inputenc}
\usepackage[margin=1in]{geometry}
\usepackage{setspace}
\doublespacing

\title{Letter to the Narrator of Alice Munro's \textit{Boys and Girls}}

\author{%
  Ma Toan Bach\\
  LSO100KCC\\
  Letter to a Character
}

\date{}

\begin{document}

\maketitle

\noindent Dear young narrator,

After reading Alice Munro's \textit{Boys and Girls}, I kept thinking about how quietly your world changes around you. At the beginning of the story, you seem so certain about who you are. You work outside with your father, help with the foxes, and take pride in being useful in a space that feels open, important, and alive. You do not see yourself as weak or secondary. You see yourself as someone who can work, observe, imagine, and matter. That is why the ending, when your father says, ``She's only a girl,'' feels so painful. It is not just a comment about one mistake. It is a sentence that tries to shrink your whole identity.

I understand why you are drawn to your father's world. The outside work gives you a sense of purpose that the house does not. Your father's work may be harsh, especially when the horses Mack and Flora are brought in to be killed for fox meat, but you still feel closer to freedom there than in the kitchen. Your mother wants you to help indoors, and she speaks as if this is natural and practical, but you experience it as a threat. I do not think you dislike your mother simply because she is your mother. I think you fear what her world represents: a future where your choices have already been made for you.

That is why Laird's growing importance must hurt. He is younger than you, but because he is a boy, the family begins to imagine him as your father's real helper. You have worked hard to earn your place, yet gender seems to matter more than effort. When your mother says that Laird will be ``a real help'' someday, it suggests that your help has always been temporary. I wonder how lonely that made you feel, especially because no one explains the change directly. They simply expect you to understand that being a girl means moving indoors, becoming quieter, and giving up the freedom you once claimed.

Your choice to open the gate for Flora is the moment I find most complicated. Part of me wants to ask why you did it, since you knew your father needed the horse and you knew there would be consequences. But I also understand that Flora's escape means more than disobedience. Flora becomes a living image of the freedom you are losing. When she runs, you recognize something powerful in her refusal to be led calmly toward death. Opening the gate seems like your own brief act of resistance, even if you cannot fully explain it afterward.

At the dinner table, when everyone looks at you and Laird tells what happened, your silence says a lot. You seem ashamed, but also trapped. Your father does not really ask who you are or what you felt. Instead, he explains you away: ``She's only a girl.'' I disagree with that judgment. You are not ``only'' anything. You are a young person beginning to understand that society can turn gender into a limit before a person has chosen it for herself.

Still, I do not read your final thought, ``Maybe it was true,'' as simple defeat. To me, it shows how powerful those expectations are. When people repeat an idea often enough, even an unfair one can start to sound believable. Your story helped me understand that growing up is not always a free discovery of identity. Sometimes it is also a painful lesson in the roles other people expect us to accept.

I hope you remember the courage in your choice at the gate. Even if you could not defend it at the table, it showed that some part of you still recognized freedom when you saw it.

\vspace{1em}

\noindent Sincerely,\\
Ma Toan Bach

\end{document}
